\uuid{qv5m}
\exo7id{1747}
\titre{exo7 1747}
\auteur{maillot}
\organisation{exo7}
\datecreate{2001-09-01}
\isIndication{false}
\isCorrection{false}
\chapitre{Topologie}
\sousChapitre{Autre}
\cpge{mpsi-mp2i.raisonnement-vocabulaire-ensembliste.c, pcsi.raisonnement-vocabulaire-ensembliste.d, ptsi.raisonnement-vocabulaire-ensembliste.d, bcpst1.vocabulaire-applications, mp-mpi.topologie-espaces-normes.d, pc.espaces-vectoriels-normes.d, psi.espaces-vectoriels-normes.d, pt.fonctions-deux-trois-variables.a}
\module{Analyse}
\niveau{L2}
\difficulte{}

\contenu{
\question{
Soit l'application
$\left(\begin{array}{cccl}
G:&\R^2&\longrightarrow&\R^2\\
&(u,v)&\longmapsto&(\frac{u}{u+v},\frac{\sqrt{v(v+2u)}}{u+v})
\end{array}\right)$. On note
$\mathcal{D}$ l'ensemble de définition de $G$. Déterminer $G(\mathcal{D})$.
}
}
